%% -------------------------------------------------------
%% Kapitel 1: Einleitung
%% -------------------------------------------------------
\chapter{Einleitung}
\label{chap:einleitung}

\section{Ausgangssituation und Problemstellung}
\label{sec:ausgangssituation}

Wer mit einem Hund in einer österreichischen Stadt oder Gemeinde unterwegs ist, sieht sich
täglich mit einer Vielzahl praktischer Fragen konfrontiert: Wo befindet sich die nächste
Freilaufwiese, auf der der Hund ohne Leine toben darf? Gibt es in der Nähe einen
Hundekot-Müllbehälter oder einen Sackerlspender? Und – besonders relevant für Halter\-innen
und Halter bestimmter Rassen – gilt an diesem Ort eine Leinen- oder Maulkorbpflicht für
ihren Hund?

Diese Informationen sind zwar grundsätzlich verfügbar, aber über viele verschiedene Quellen
verstreut: kommunale Webseiten, amtliche Kundmachungen, geografische Open-Data-Portale oder
mündliche Weitergabe unter Hundehalter\-innen und Hundehaltern. Eine zentrale, standortbewusste
und rassenspezifische Anlaufstelle fehlt im deutschsprachigen App-Markt weitgehend.

\subsection{Defizite bestehender Lösungen}
\label{subsec:defizite}

Eine Analyse verfügbarer Anwendungen zeigt, dass keine der bekannten Lösungen alle
relevanten Anforderungen in einer einzigen App vereint:

\begin{table}[H]
  \centering
  \caption{Vergleich bestehender Lösungen}
  \label{tab:konkurrenzanalyse}
  \small
  \begin{tabularx}{\textwidth}{|p{2.8cm}|Y|Y|Y|Y|}
    \hline
    \textbf{App / Dienst} & \textbf{Karte mit POIs} & \textbf{Rechtsvorgaben} &
    \textbf{Rassenfilter} & \textbf{iOS-App} \\
    \hline
    BringGassi       & Ja (wenige) & Nein & Nein & Ja \\
    \hline
    Google Maps      & Nein        & Nein & Nein & Ja \\
    \hline
    Open\-Street\-Map & Begrenzt   & Nein & Nein & Indirekt \\
    \hline
    Wien.gv.at       & Ja (Wien)   & Begrenzt & Nein & Nein \\
    \hline
    \textbf{DoggoGO} & \textbf{Ja} & \textbf{Ja} & \textbf{Ja} & \textbf{Ja} \\
    \hline
  \end{tabularx}
\end{table}

\textbf{BringGassi} bietet zwar eine kartenzentrierte Ansicht für Hundeausflüge, deckt
jedoch keine rechtlichen Vorgaben ab und ist regional stark begrenzt. \textbf{Google Maps}
erlaubt grundsätzlich die Suche nach Hundeauslaufzonen, verfügt aber über keinerlei
rassenspezifische Rechtsinformationen und ist nicht auf Hundehaltung spezialisiert. Kommunale
Portale wie das der Stadt Wien bieten Standortdaten für ihren Zuständigkeitsbereich, sind
aber weder mobil optimiert noch rassenspezifisch und nicht überregional einsetzbar.

Das Kerndefizit aller betrachteten Lösungen liegt in der fehlenden Verknüpfung von drei
Aspekten, die Hundehalter\-innen und Hundehalter täglich benötigen:
\begin{enumerate}
  \item standortbezogene Karte mit relevanten Points of Interest (POIs),
  \item rassenspezifische gesetzliche Informationen (Leinen-/Maulkorbpflicht),
  \item plattformübergreifende Verfügbarkeit (Web \emph{und} iOS).
\end{enumerate}

\section{Zielsetzung dieser Arbeit}
\label{sec:zielsetzung}

Diese Diplomarbeit behandelt die Realisierung von zwei zentralen Komponenten des
DoggoGO-Projekts: dem \textbf{Backend} und der \textbf{iOS-Applikation}.

Das \textbf{Backend} hat folgende Aufgaben:
\begin{itemize}
  \item Bereitstellung einer RESTful~API für die Angular-Webanwendung und die iOS-App,
  \item Verwaltung von Hundehalter\-innen und Hundehaltern sowie ihrer Hundeprofile,
  \item Persistenz der Hunderassen inklusive ihrer rassenspezifischen Rechtsvorschriften,
  \item sicherer Betrieb in einer Docker-Container-Umgebung.
\end{itemize}

Die \textbf{iOS-App} soll:
\begin{itemize}
  \item eine interaktive Karte mit standortbasierten Markierungen (Freilaufwiesen,
        Sackerlspender, Hundekot-Müllbehälter) anzeigen,
  \item das Hundeprofil der Nutzerin / des Nutzers verwalten,
  \item rassenspezifische Leinen- und Maulkorbvorschriften für den aktuellen Standort anzeigen,
  \item eine sichere Anmeldung per JWT-Authentifizierung bieten.
\end{itemize}

\section{Abgrenzung}
\label{sec:abgrenzung}

Die Angular-Webanwendung (inklusive UI/UX-Konzept, Routenplanung, Wetterintegration und
Progressive-Web-App-Funktionalität) wurde von Manuel Rankl entwickelt und ist nicht
Gegenstand dieser Arbeit. Gleiches gilt für die Entwicklung eines Android-Clients. Das in
dieser Arbeit beschriebene Backend bildet die gemeinsame Datenbasis für beide Frontends.

\section{Aufbau der Arbeit}
\label{sec:aufbau}

Die vorliegende Arbeit ist wie folgt gegliedert:

\begin{description}
  \item[Kapitel~\ref{chap:anforderungen} – Anforderungsanalyse:]
    Beschreibung der Stakeholder, funktionaler und nicht-funktionaler Anforderungen sowie
    der Use-Case-Szenarien aus Sicht des Backends und der iOS-App.
  \item[Kapitel~\ref{chap:architektur} – Systemarchitektur:]
    Gesamtarchitektur der Plattform, Datenbankdesign, REST-API-Überblick und
    Kommunikationsfluss zwischen den Systemkomponenten.
  \item[Kapitel~\ref{chap:technologie} – Technologieauswahl:]
    Begründete Auswahl der eingesetzten Technologien (ASP.NET~Core~8, Entity~Framework~Core~9,
    PostgreSQL, iOS/SwiftUI) mit Bewertung von Alternativen.
  \item[Kapitel~\ref{chap:backend} – Backend-Implementierung:]
    Detaillierte Beschreibung der API-Endpunkte, Datenbankmodelle, Migrations,
    Authentifizierung, Breed-Seeding und Docker-Deployment.
  \item[Kapitel~\ref{chap:ios} – iOS-Applikation:]
    Architektur und Implementierung der iOS-App: MapKit-Integration, Standortdienste,
    Authentifizierungsflow und Hundeprofilverwaltung.
  \item[Kapitel~\ref{chap:testing} – Testing und Deployment:]
    Teststrategie, durchgeführte Tests und Erkenntnisse aus dem Deployment.
  \item[Kapitel~\ref{chap:zusammenfassung} – Zusammenfassung:]
    Reflexion der erreichten Ziele, suboptimaler Entscheidungen, Einschränkungen und
    Erweiterungsmöglichkeiten.
\end{description}
